% !TeX root = ./../../main.tex
% Rysunek fig:cls:pdu-curves: dokładność modelu M2 (c7_f200_q3_b13_d2) w funkcji siły jednej składowej szumu,
% bez korekty odczytu i z korektą w wariantach FZ i AR (eksperyment symulacyjny, 2000 próbek testowych).
% Wykres odtworzony w pgfplots z plików danych; żadna wartość nie została odczytana z rysunków.
% Dane: figures/chp_classification/cls_pdu_m2_<składowa>_<pomiary>.csv, czyli niezmienione kopie plików
%   eksperymenty/krzywe-degradacji/klasyfikator/wyniki/pgfplots/c7_f200_q3_b13_d2__<składowa>__<pomiary>.csv
%   (składowa: odstrojenie, tlumienie_fazy, depolaryzacja; pomiary: 1024 albo dokladne).
% Kolumny: sila, dokl_bez (bez korekty), dokl_a2 (wariant FZ), dokl_b2 (wariant AR); pierwszy wiersz (sila = 0)
%   to brak szumu obwodu i jest pomijany, bo oś jest logarytmiczna.
% Poziom ibm_fez modelu M2: odstrojenie 0,0012756 rad, depolaryzacja 0,0029837, tłumienie fazy 0 (kolumna ibm_fez = 1).
% Najlepsza stała predykcja: 0,1555 (wyniki/kontrole.csv, kolumna dokl_stala_predykcja).
% Wymaga: pgfplots z biblioteką groupplots i biblioteki calc (wczytane w config/my-settings.tex).
\begin{tikzpicture}
\pgfplotsset{
  clsbez/.style={black, thick, mark=*, mark size=1.5pt},
  clsexact/.style={black!45, line width=1.3pt, no marks},
  clsfz/.style={black, dashed, mark=square, mark size=1.6pt, mark options={solid}},
  clsar/.style={black, densely dashdotted, mark=triangle, mark size=2.2pt, mark options={solid}},
}
\begin{groupplot}[
  group style={group size=3 by 1, horizontal sep=0.45cm, y descriptions at=edge left},
  width=0.385\textwidth,
  height=4.5cm,
  xmode=log,
  xmin=0.001, xmax=0.7,
  xtick={0.001,0.01,0.1},
  ymin=0.05, ymax=0.47,
  ytick={0.1,0.2,0.3,0.4},
  ymajorgrids,
  grid style={black!15},
  tick label style={font=\footnotesize},
  label style={font=\footnotesize},
  title style={font=\footnotesize, yshift=-1mm},
  yticklabel style={/pgf/number format/.cd, fixed, fixed zerofill, precision=1, use comma},
  legend style={font=\scriptsize, draw=black!40, fill=white, inner sep=2pt, /tikz/every even column/.append style={column sep=3mm}},
  legend cell align=left,
  legend columns=4,
]
% ---------- (a) odstrojenie ----------
\nextgroupplot[
  title={(a) odstrojenie},
  xlabel={$\cohangle$ [rad na blok]},
  ylabel={dokładność},
  legend to name=clspdulegend,
]
\draw[black!40, thin, densely dotted] (axis cs:0.001,0.1555) -- (axis cs:0.7,0.1555);
\draw[black!55, thin, densely dashed] (axis cs:0.0012756,0.05) -- (axis cs:0.0012756,0.47);
\node[font=\scriptsize, rotate=90, anchor=south west, inner sep=1.5pt] at (axis cs:0.0021,0.06) {ibm\_fez};
\addplot[clsexact, skip coords between index={0}{1}] table[col sep=comma, x=sila, y=dokl_bez] {figures/chp_classification/cls_pdu_m2_odstrojenie_dokladne.csv};
\addlegendentry{bez korekty, wartości dokładne}
\addplot[clsbez, skip coords between index={0}{1}] table[col sep=comma, x=sila, y=dokl_bez] {figures/chp_classification/cls_pdu_m2_odstrojenie_1024.csv};
\addlegendentry{bez korekty}
\addplot[clsfz, skip coords between index={0}{1}] table[col sep=comma, x=sila, y=dokl_a2] {figures/chp_classification/cls_pdu_m2_odstrojenie_1024.csv};
\addlegendentry{wariant FZ}
\addplot[clsar, skip coords between index={0}{1}] table[col sep=comma, x=sila, y=dokl_b2] {figures/chp_classification/cls_pdu_m2_odstrojenie_1024.csv};
\addlegendentry{wariant AR}
% ---------- (b) tłumienie fazy ----------
\nextgroupplot[
  title={(b) tłumienie fazy},
  xlabel={$\qphase$},
]
\draw[black!40, thin, densely dotted] (axis cs:0.001,0.1555) -- (axis cs:0.7,0.1555);
\addplot[clsexact, skip coords between index={0}{1}] table[col sep=comma, x=sila, y=dokl_bez] {figures/chp_classification/cls_pdu_m2_tlumienie_fazy_dokladne.csv};
\addplot[clsbez, skip coords between index={0}{1}] table[col sep=comma, x=sila, y=dokl_bez] {figures/chp_classification/cls_pdu_m2_tlumienie_fazy_1024.csv};
\addplot[clsfz, skip coords between index={0}{1}] table[col sep=comma, x=sila, y=dokl_a2] {figures/chp_classification/cls_pdu_m2_tlumienie_fazy_1024.csv};
\addplot[clsar, skip coords between index={0}{1}] table[col sep=comma, x=sila, y=dokl_b2] {figures/chp_classification/cls_pdu_m2_tlumienie_fazy_1024.csv};
% ---------- (c) depolaryzacja ----------
\nextgroupplot[
  title={(c) depolaryzacja},
  xlabel={$\qdep$},
]
\draw[black!40, thin, densely dotted] (axis cs:0.001,0.1555) -- (axis cs:0.7,0.1555);
\draw[black!55, thin, densely dashed] (axis cs:0.0029837,0.05) -- (axis cs:0.0029837,0.47);
\node[font=\scriptsize, rotate=90, anchor=south west, inner sep=1.5pt] at (axis cs:0.0048,0.06) {ibm\_fez};
\addplot[clsexact, skip coords between index={0}{1}] table[col sep=comma, x=sila, y=dokl_bez] {figures/chp_classification/cls_pdu_m2_depolaryzacja_dokladne.csv};
\addplot[clsbez, skip coords between index={0}{1}] table[col sep=comma, x=sila, y=dokl_bez] {figures/chp_classification/cls_pdu_m2_depolaryzacja_1024.csv};
\addplot[clsfz, skip coords between index={0}{1}] table[col sep=comma, x=sila, y=dokl_a2] {figures/chp_classification/cls_pdu_m2_depolaryzacja_1024.csv};
\addplot[clsar, skip coords between index={0}{1}] table[col sep=comma, x=sila, y=dokl_b2] {figures/chp_classification/cls_pdu_m2_depolaryzacja_1024.csv};
\end{groupplot}
\node[anchor=north] at ($(group c2r1.south)+(0,-1.05cm)$) {\pgfplotslegendfromname{clspdulegend}};
\end{tikzpicture}
